%! AP Statistics — Unit 1, Lesson 1 — Lesson Plan
\documentclass[10pt]{article}
\usepackage[margin=0.6in]{geometry}
\usepackage{apstats-colors}
\usepackage[most]{tcolorbox}
\usepackage{enumitem}
\usepackage{multicol}
\usepackage{amsmath,amssymb}
\usepackage{array}
\usepackage{tabularx}
\usepackage{apstats-article}
\usepackage{apstats-boxes}
\usepackage{apstats-colors}
\tcbuselibrary{breakable}

\newcommand{\UnitNumberName}{Unit 1: Exploring One-Variable Data \quad}
\newcommand{\LessonNumberName}{Lesson 1.1: Introducing Statistics --- What Can We Learn from Data?}

\begin{document}
\begin{center}
    {\Large\bfseries \CourseName: \SchoolYear\ (\MeetingLength) \\
    {\small\bfseries \UnitNumberName \LessonNumberName}}
\end{center}

\begin{tcolorbox}[colback=sky, colframe=navy, boxrule=0.9pt, arc=2mm,
    left=2mm, right=2mm, top=1.5mm, bottom=1.5mm]
    \textbf{Primary Objective:} Students will understand that statistics is the science of
    collecting, analyzing, and drawing conclusions from data. Students will distinguish between
    populations and samples, identify individuals and variables in a dataset, and recognize why
    variation in data is central to statistical thinking (Big Idea: VAR).
\end{tcolorbox}

\renewcommand{\arraystretch}{1.6}
\begin{skillbox}[Priority Ideas \& Skills]{goldbox}
\begin{minipage}[t]{0.48\linewidth}
    \textbf{AP Skill 1 --- Selecting Statistical Methods}
    \begin{itemize}
        \item Identify the population of interest and the sample in a given scenario.
        \item Explain why a sample is used instead of a census.
        \item Recognize that statistics answers questions using data collected with variability.
    \end{itemize}
\end{minipage}
\hfill
\begin{minipage}[t]{0.48\linewidth}
    \textbf{Key Understandings}
    \begin{itemize}
        \item Statistics $\neq$ a single number; it is a process of inquiry.
        \item Variation is natural and expected; statistical methods quantify it.
        \item The same question can yield different answers in different samples --- and that is informative, not a mistake.
    \end{itemize}
\end{minipage}
\end{skillbox}

\begin{skillbox}[{Vocabulary, Concepts, \& Theorems}]{greenbox}
\begin{minipage}[t]{0.48\linewidth}
    \begin{itemize}
        \item \textbf{Statistics} --- science of collecting, analyzing, and drawing conclusions from data
        \item \textbf{Individual} --- a single person, animal, object, or event in a dataset
        \item \textbf{Population} --- the entire group of individuals of interest
        \item \textbf{Sample} --- a subset of the population from which data are actually collected
    \end{itemize}
\end{minipage}
\hfill
\begin{minipage}[t]{0.48\linewidth}
    \begin{itemize}
        \item \textbf{Census} --- data collected from every individual in a population
        \item \textbf{Variability / Variation} --- natural differences in values across individuals
        \item \textbf{Data} --- recorded values for individuals; can be categorical or quantitative
        \item \textbf{Anecdotal evidence} --- data from a few unrepresentative cases; not reliable for inference
    \end{itemize}
\end{minipage}
\end{skillbox}

\begin{skillbox}[Activate Prior Knowledge \& Spiral Review]{sky}
\begin{minipage}[t]{0.48\linewidth}
    \textbf{Prior Knowledge Check (3--4 min)}
    \begin{itemize}
        \item Ask: \textit{``When you hear the word `average,' what does that tell you? What doesn't it tell you?''}
        \item Have students share one question they think data could answer about their school (e.g., ``How long do students sleep?'').
    \end{itemize}
\end{minipage}
\hfill
\begin{minipage}[t]{0.48\linewidth}
    \textbf{Connections to Future Units}
    \begin{itemize}
        \item The idea of a \textit{sample} vs.\ \textit{population} is foundational for all inference in Units 6--9.
        \item Variability in data motivates the need for summary statistics (Lessons 1.6--1.7) and probability (Unit 4).
    \end{itemize}
\end{minipage}
\end{skillbox}

\begin{skillbox}[Hook \& Lesson]{sky}
\begin{minipage}[t]{0.48\linewidth}
    \textbf{Hook (5--7 min)}\\
    Display a headline with a surprising statistic (e.g., ``Students who sleep more score higher on tests''). Ask:
    \begin{itemize}
        \item How do you think they know this?
        \item Did they study every student in the world?
        \item Could this be wrong?
    \end{itemize}
    Transition: \textit{``These are exactly the questions statistics is designed to answer.''}
\end{minipage}
\hfill
\begin{minipage}[t]{0.48\linewidth}
    \textbf{Lesson Flow (15--18 min)}
    \begin{enumerate}
        \item Introduce the \textit{statistical problem-solving process}: Ask a question $\to$ collect data $\to$ analyze $\to$ interpret.
        \item Use a class example: ``What is a typical commute time for adults in our city?'' --- identify the individual, population, and sample.
        \item Show a small dataset (5--6 rows); identify individuals, variables, and values.
        \item Discuss: Why can't we always do a census? (Cost, time, impossibility.)
    \end{enumerate}
\end{minipage}
\end{skillbox}

\begin{skillbox}[Explicit Instruction \& Active Monitoring]{sky}
\begin{minipage}[t]{0.48\linewidth}
    \textbf{Direct Instruction (10 min)}
    \begin{itemize}
        \item Walk through 2--3 real contexts (sports, medicine, social media), each time asking students to identify: individual, population, sample, and at least one variable.
        \item Emphasize: variation is the \textit{reason} statistics is interesting --- if everyone were identical, we would not need statistics.
    \end{itemize}
\end{minipage}
\hfill
\begin{minipage}[t]{0.48\linewidth}
    \textbf{Active Monitoring}
    \begin{itemize}
        \item Circulate while students attempt a quick identification task (see Group Work).
        \item Listen for common misconception: confusing \textit{sample} with \textit{data} (a sample is the set of individuals; the data are the recorded values).
        \item Cold-call 2--3 students to share their reasoning.
    \end{itemize}
\end{minipage}
\end{skillbox}

\begin{skillbox}[Group Work \& Differentiation]{redbox}
\begin{minipage}[t]{0.48\linewidth}
    \textbf{Partner Activity (8--10 min)}\\
    Provide 3 statistical scenarios (e.g., a poll of 500 voters, a study of 30 mice, a school survey). For each, groups identify and record:
    \begin{itemize}
        \item The research question
        \item The individual
        \item The population
        \item The sample (if applicable)
        \item One variable that was likely measured
    \end{itemize}
\end{minipage}
\hfill
\begin{minipage}[t]{0.48\linewidth}
    \textbf{Differentiation}
    \begin{itemize}
        \item \textit{Support}: Provide a sentence frame: ``The individuals are \underline{\hspace{1cm}}. The population is \underline{\hspace{1cm}}.''
        \item \textit{Extension}: Ask students to write their \textit{own} statistical question about the class, identify all components, and explain why a census is or is not feasible.
        \item \textit{ELL}: Pair vocabulary list with visual icons (e.g., one person = individual; crowd = population).
    \end{itemize}
\end{minipage}
\end{skillbox}

\begin{skillbox}[Individual Work \& Assessment]{redbox}
\begin{minipage}[t]{0.48\linewidth}
    \textbf{Exit Ticket (5 min)}\\
    Present one new scenario (e.g., a hospital tracks recovery times for 50 patients after surgery). Students independently write:
    \begin{enumerate}
        \item The individual
        \item The population
        \item The sample
        \item One variable and whether it varies
    \end{enumerate}
\end{minipage}
\hfill
\begin{minipage}[t]{0.48\linewidth}
    \textbf{AP-Style Check}\\
    \textit{Multiple-choice practice:} ``A researcher surveys 200 randomly chosen adults about their exercise habits. The population is\ldots'' \\[4pt]
    Use student responses to gauge readiness for Lesson 1.2 (types of variables). Collect exit tickets to identify gaps.
\end{minipage}
\end{skillbox}

\begin{skillbox}[Reinforcement \& Extension]{goldbox}
\begin{minipage}[t]{0.48\linewidth}
    \textbf{Homework / Practice}
    \begin{itemize}
        \item Read 3 short news blurbs with statistics; for each, identify the individual, population, sample (if stated), and one variable.
        \item Reflect: \textit{``Why might the results of this study not apply to everyone?''}
    \end{itemize}
\end{minipage}
\hfill
\begin{minipage}[t]{0.48\linewidth}
    \textbf{Extension / Preview}
    \begin{itemize}
        \item \textit{Extension}: Research the U.S.\ Census --- what are the costs and challenges of conducting one every 10 years?
        \item \textit{Preview}: Lesson 1.2 will classify variables as \textit{categorical} or \textit{quantitative} --- start noticing what type of answer each variable you identified today would produce.
    \end{itemize}
\end{minipage}
\end{skillbox}

\end{document}
